%% Chapter 3, Section 3.5: Exercises
%% A First Course in Monte Carlo Methods — Sanz-Alonso & Al-Ghattas

\section{Exercises}
\label{sec:3.5}

\begin{enumerate}

\item[3.1] Let $f_\nu$ denote the p.d.f.\ of Student's $t$-distribution with
$\nu$ degrees of freedom. Using a sample size $N = 1000$, estimate
$\Expect_{X \sim f_5}[|X|]$ by:
\begin{enumerate}
  \item[(a)] Classical Monte Carlo integration.
  \item[(b)] Importance sampling with reference to the $f_1$ distribution.
  \item[(c)] Importance sampling with reference to the $\Normal(0,1)$
             distribution.
\end{enumerate}
Compare the accuracy of the different methods by estimating the variance of the
estimators, replicating the experiment $M = 10^4$ times.

\item[3.2] In this exercise, you will learn how to modify the autonormalized
importance sampling algorithm to compute weights in the log-scale, which is
helpful to avoid overflow or underflow. First, define
$w_{\log}(x) = \log f(x) - \log g(x)$. Next, given a sample
$X^{(1)}, \ldots, X^{(N)} \iid g$, let
$w_{\max} = \max_{1 \leq n \leq N} w_{\log}(X^{(n)})$, and let
$w^{(n)} = \exp\!\bigl(w_{\log}(X^{(n)}) - w_{\max}\bigr)$. Finally, normalize
the weights by setting
$w^*(X^{(n)}) := w^{(n)} / \sum_{l=1}^{N} w^{(l)}$. Implement an
autonormalized importance sampling algorithm using this calculation of the
weights in the log-scale to estimate the expected value of a $\Normal(0,1)$
using as proposal a $\Normal(1,2)$. Use a sample size $N = 10^5$.

\item[3.3] In the setting of Example~\ref{ex:3.7}, choose $h(x) = (1+x)^{-1}$.
Approximate the value of the integral (which is $\log 2$) using a sample size
$N = 1500$ with classical Monte Carlo and antithetic sampling. Compute the
variance of the estimators by repeating the experiment a total of $M = 10^4$
times. Does antithetic sampling give variance reduction?

\item[3.4] In the setting of Example~\ref{ex:3.8}, choose $h(x) = e^x$.
Approximate the value of the integral (which is $e - 1$) using a sample size
$N = 1500$ with classical Monte Carlo and control variates. For the control
variates implementation, choose $\hat{h}(x) = 1 + x + x^2/2 + x^3/6$ and find
$\mu$ analytically. Compute the variance of the estimators by repeating the
experiment a total of $M = 10^4$ times. Do you achieve variance reduction via
control variates?

\item[3.5] In this exercise, you will learn another variance reduction
technique: stratified sampling. Let $h : [0,1] \to \R$ be such that
$\int_0^1 h^2(x)\,dx < \infty$ and consider the following approximations of
$\calI[h] := \int_0^1 h(x)\,dx$:
\[
  \calI^{\mathrm{mc}}[h]
  = \frac{1}{N}\sum_{n=1}^{N} h(U^{(n)}),
  \qquad U^{(n)} \iid \Unif(0,1),
\]
\[
  \calI^{\mathrm{ss}}[h]
  = \frac{1}{N}\sum_{n=1}^{N} h(V^{(n)}),
  \qquad V^{(n)} \sim \Unif\!\bigl((n-1)/N,\, n/N\bigr),
\]
where the $V^{(n)}$ are sampled independently.
\begin{enumerate}
  \item[(a)] Show that the \emph{stratified sampling} estimator
             $\calI^{\mathrm{ss}}[h]$ is unbiased and that
             \[
               \Var\!\left[\calI^{\mathrm{ss}}[h]\right]
               \leq \Var\!\left[\calI^{\mathrm{mc}}[h]\right].
             \]
  \item[(b)] Show that if $h$ is Lipschitz, then the stratified sampling
             estimator has a variance that decays as $N^{-3}$.
\end{enumerate}

\item[3.6] The p.d.f.\ $f_\theta(x)$ is in the exponential family
$\mathcal{E}_F$ if it can be written in the form
$f_\theta(x) = e^{\langle t(x),\theta\rangle - F(\theta) + k(x)}$. We denote
this distribution as $\mathcal{E}_F(\theta)$. Suppose the natural parameter
space $\Theta = \{\theta \mid \int_{\R^d} f_\theta(x)\,dx = 1\}$ is affine,
meaning that $\sum_i w_i\theta_i \in \Theta$ if $\theta_i \in \Theta$ and
$\sum_i w_i = 1$. Show that the $\chi^2$-divergence within the same exponential
family $\mathcal{E}_F$ is characterized by
\[
  d_{\chi^2}\!\bigl(\mathcal{E}_F(\theta_1)\|\mathcal{E}_F(\theta_2)\bigr)
  = e^{F(2\theta_1 - \theta_2) - 2F(\theta_1) + F(\theta_2)} - 1.
\]

\item[3.7] Let $f = \Normal(\mu, \Sigma)$ be a multivariate normal
distribution.
\begin{enumerate}
  \item[(a)] Show that it belongs to the exponential family, with parameter
             $\theta = [\Sigma^{-1}\mu;\, -\tfrac{1}{2}\Sigma^{-1}]$ where
             the natural parameter space $\R^d \oplus \R^{d \times d}$ inherits
             the inner products from the Euclidean space and the matrix
             space,\footnote{The canonical inner product in the space of square
             matrices is defined to be the trace of the matrix product. We can
             also view this as an extension of the Euclidean inner product where
             the scalar in each coordinate is replaced by vectors.}
             $t(x) = [x;\, xx^T]$,
             $F(\theta) = \tfrac{1}{2}\mu^T\Sigma^{-1}\mu
               + \tfrac{1}{2}\log|\Sigma|$, and
             $k(x) = -\tfrac{d}{2}\log(2\pi)$.

  \item[(b)] Let $f_1 = \Normal(\mu_1, \Sigma_1)$ and
             $f_2 = \Normal(\mu_2, \Sigma_2)$ be two multivariate normal
             distributions. Show that
             \[
               d_{\chi^2}(f_1\|f_2)
               = \frac{|W|}{\sqrt{|2W - I|}}\,
                 e^{u^T \Sigma^{-1}(2W-I)^{-1} u}
                 - 1,
             \]
             where $W = \Sigma_2 \Sigma_1^{-1}$ and $u = \mu_1 - \mu_2$.
\end{enumerate}

\end{enumerate}
